% ---------------------------------------------------------------------
% Drop-in replacement for the main-text passage (98)-(103) of Sec. VI-E:
% how the identified onset responds to a perturbation.
%
% What this changes against the current draft:
%   - (99) introduces the true pre-onset level b, so that C has something
%     to estimate; without it C appears in (98) and vanishes by (102)
%     with no accounting, which reads as a dropped term.
%   - a sentence fixing e_omega to the TRUE onset, so that it carries no
%     dependence on the optimisation variable.
%   - the Leibniz boundary terms are named and shown to cancel
%     identically, rather than being passed over.  This is the first
%     thing a reader checks when a window endpoint is also the variable
%     being optimised.
%   - (104)-(105) carry the baseline error term the earlier flow lost,
%     and bound it in closed form.
%   - the measured sizes are NOT here: they belong to the realised
%     budget, Sec. VIII-E.  Only the a priori gain range is quoted.
%
% Symbols follow the manuscript: t_0 record start, t_c onset, t_c^* its
% true value, \hat t_c the minimiser, x = C_2 tau_end.
%
% Equation numbers continue from (97); renumber if the host document
% differs.
% ---------------------------------------------------------------------

\documentclass[10pt]{article}
\usepackage[a4paper,margin=22mm]{geometry}
\usepackage{amsmath,amssymb,booktabs}
\setlength{\parskip}{0.35em}\setlength{\parindent}{0pt}
\newcommand{\Wz}{W\!z_{CoM}}
\newcommand{\tc}{t_{\rm c}}
\newcommand{\tcs}{t_{\rm c}^{*}}
\newcommand{\tend}{t_{\rm end}}
\title{\vspace{-16mm}\textbf{Sec.~VI-E, (98)--(105): the onset under a
perturbation}\vspace{-3mm}}
\date{\small drop-in main text; the measured sizes live in Sec.~VIII-E}
\begin{document}\maketitle

Because the measured angular rate $y(t)$ can be regarded as the nominal
response perturbed by the unmodelled forcing $\rho$, the minimiser
selects $\hat\tc$ by minimising the two-segment cost over the record.
Both segments enter, and both share the pre-onset level $C$:
\begin{equation}
\begin{aligned}
  J(\tc) = &\int_{t_0}^{\tc}\bigl(y(t) - C\bigr)^2dt \\
    &+ \int_{\tc}^{\tend}
       \Bigl(y(t) - \bigl(\hat\omega(t;\tc) + C\bigr)\Bigr)^2dt .
\end{aligned}
  \tag{98}
\end{equation}

Write the recorded rate as the true pre-onset level, the nominal
response from the true onset, and a deviation:
\begin{equation}
  y(t) = b + \hat\omega(t;\tcs) + e_\omega(t),
  \qquad
  \hat\tc = \tcs + \delta \tc .
  \tag{99}
\end{equation}
Here $b$ is the rate the vehicle actually holds before the excitation,
zero under the onset conditions. The rate is taken from the navigation
filter, whose delta-angle bias states remove the gyro bias, so $b$
represents residual pre-excitation motion rather than sensor bias; $C$
is its estimate, formed over the pre-onset segment. Note also that
$e_\omega$ is anchored at the \emph{true} onset $\tcs$: once a run is
recorded it is a fixed function of $t$, and it carries no dependence on
the optimisation variable.

A first-order expansion of the estimated rate gives
\begin{equation}
  \hat\omega\bigl(t;\hat\tc\bigr)
  \;\simeq\; \hat\omega\bigl(t;\tcs\bigr) + \delta\tc\,\chi(t)
  \tag{100}
\end{equation}
with the sensitivity
\begin{equation}
  \chi(t) \triangleq
  \frac{\partial\hat\omega(t;\tc)}{\partial \tc}\bigg|_{\tc = \tcs}
  = -\,C_1C_2\sinh C_2\bigl(t - \tcs\bigr).
  \tag{101}
\end{equation}

Two features of (98) make the differentiation clean, and both are worth
stating because each looks like an omission otherwise.

First, $\tc$ is the upper limit of one integral and the lower limit of
the other, so the Leibniz rule produces two boundary terms. They cancel
identically. The rise contributes $C_1(\cosh 0 - 1) = 0$ at $t = \tc$,
so the two segments agree there and both terms equal
$\bigl(y(\tc) - C\bigr)^2$ with opposite signs. Nothing is discarded,
and the cancellation is exact rather than asymptotic.

Second, substituting (99) and (100), the post-onset residual is
\begin{equation}
  y(t) - \hat\omega(t;\hat\tc) - C
  = e_\omega(t) - \delta\tc\,\chi(t) - (C - b),
  \tag{102}
\end{equation}
so $C$ enters only through the baseline \emph{error} $C - b$, not
through its own size. A large but correctly estimated pre-onset level
is harmless.

Differentiating (98) with respect to $\delta\tc$ --- the same operator
as $\partial/\partial\tc$ by (99), but the variable in which the
linearised cost is an exact quadratic --- and setting the result to zero
gives, when the pre-onset fit is unbiased,
\begin{equation}
  \frac{\partial J}{\partial\,\delta\tc}
  = -2\int_{\tc}^{\tend}\bigl(e_\omega(t) - \delta\tc\,\chi(t)\bigr)
    \chi(t)\,dt = 0 ,
  \tag{103}
\end{equation}
which implies
\begin{equation}
  \delta\tc = \frac{\langle e_\omega,\chi\rangle}{\|\chi\|^{2}} .
  \tag{104}
\end{equation}
Here $\langle a,b\rangle$ denotes the $L^2$ inner product over the
fitting window $[\tcs,\tend]$, i.e.\ $\langle a,b\rangle \triangleq
\int_{\tcs}^{\tend}a(t)b(t)\,dt$, and $\|b\|^2 \triangleq \langle
b,b\rangle$. The shift $\delta\tc$ is well defined because the
maximum-angle event occurs after the onset, $\tau_{\rm end} > 0$, so
$\sinh C_2\tau$ does not vanish identically and $\|\chi\|^2 > 0$.

\paragraph{The baseline error, retained.}
Keeping the last term of (102) rather than assuming it away replaces
(104) by
\begin{equation}
  \delta\tc = \frac{\langle e_\omega,\chi\rangle}{\|\chi\|^{2}}
   \;-\; (C - b)\,\frac{\langle 1,\chi\rangle}{\|\chi\|^{2}} ,
  \tag{105}
\end{equation}
and the second term does not vanish: $\chi$ is of one sign, so
$\langle1,\chi\rangle = -C_1(\cosh x - 1)$. Evaluating it with
$\|\chi\|^2 = C_1^2C_2\bigl[\tfrac14\sinh 2x - \tfrac12 x\bigr]$ and
carrying it through $\Delta M_{\rm crit} = \dot M\,\delta\tc$, the
amplitude $C_1$ cancels as it does elsewhere and
\begin{equation}
  \Delta M_{{\rm crit},C}
  = \Wz\,(C - b)\,
    \frac{\cosh x - 1}
         {C_2\bigl[\tfrac14\sinh 2x - \tfrac12 x\bigr]} .
  \tag{106}
\end{equation}
Over the operating box of Sec.~VI-D the gain of (106) runs from
$0.042$~N\,m per rad/s at the slowest ramp to $0.351$ at the fastest,
the fast end being worse because the window is shorter. What the fitted
baselines actually are, and how far the identified offset moves when
$C$ is pinned to zero instead of fitted, are measurements and are
reported with the realised budget in Sec.~VIII-E.

\end{document}
